\subsection{Borel flags for determinant extension spaces}
\begin{finding}{P-17}{a finite Borel-flag obstruction for eighteen determinant ten-planes}
Over an algebraically closed field of characteristic zero, each
determinant residual tensor for a subset in the following table, or
its transpose, has border rank at least eighteen. There are eighteen
distinct subsets in this list and its transposes.
\[
\begin{array}{c|l|r|r|r}
 &S&\dim B&\text{quotient weights}&\text{rank exclusions}\\\hline
 1&0,1,2,3,6,7,8,12,13,14&16&101&3\\
 2&0,1,2,3,4,6,7,8,12,13&15&97&18\\
 3&0,1,2,3,6,7,8,9,12,13&15&95&19\\
 4&0,1,2,3,6,7,8,12,13,18&15&96&25\\
 5&0,1,2,3,4,6,7,8,9,12&14&90&84\\
 6&0,1,2,3,4,6,7,8,12,18&14&90&97\\
 7&0,1,2,3,4,6,7,9,12,13&14&91&61\\
 8&0,1,2,3,4,6,7,12,13,18&14&92&74\\
 9&0,1,2,3,6,7,9,12,13,15&14&88&41\\
 10&0,1,2,3,4,6,7,12,18,24&13&88&263
\end{array}
\]
This is a statement about the listed ten-planes. It does not assert
the uniform minimum $\mu_{10}\ge18$.
\end{finding}
\begin{proof}
We first prove the reduction behind the finite calculation.
Let $S$ be a coordinate order ideal in the two row/column two-subset
posets, let $W$ be the corresponding determinant quadratic space, and
let $A=W^*$. Write $B=C$ for its used-variable space, of dimension $m$,
and $T\in A\otimes B\otimes C$ for the residual inclusion tensor.
The product $H$ of the two upper triangular groups on the row and
column index spaces preserves $W$. It also preserves the used-variable
space: lowering a row or column of a used variable either stays in
its rectangle or belongs to the corresponding lower rectangle of $S$.
The contragredient action on $A$ therefore fixes $T$ and preserves
$U=T(C^*)\subset A\otimes B$.

If $\BR(T)\le r$, the proper-incidence necessary condition of R-04
gives an $r$-plane $E\supset U$ with
\[
 \rank\Psi_A(E)\le(\dim A-1)r,\qquad
 \rank\Psi_B(E)\le(m-1)r.
 \tag{P17.1}
\]
This incidence fiber is projective and $H$-invariant. The triangular
Segre--Pl\"ucker fixed-point argument given in P-11 consequently lets
us take $E$ to be $H$-invariant. Put $V=(A\otimes B)/U$ and
$F=E/U$. An $H$-stable complete flag on the triangular ambient
representation descends, after removing repeated terms, to one on
$V$. Intersect it with $F$. Since each ambient step has dimension one,
the intersection dimensions increase by at most one, and deletion of
repeated terms gives an $H$-stable complete flag on $F$.

Thus every possible $F$ is obtainable by successive invariant lines
in quotients $V/F_0$, starting with $F_0=0$. Each such line has one
torus character. The six adjacent row/column root operators must kill
it in the quotient, because a unipotent group acts trivially on a
one-dimensional representation. In characteristic zero these six
operators generate the relevant unipotent action. For the necessity
of the enumeration it already suffices that they kill every line in
an invariant flag.

For each torus weight $w$, compute the common kernel of the six
induced operators on $(V/F_0)_w$. A zero-dimensional kernel offers no
next line. A one-dimensional kernel offers exactly one, with its full
rational coordinate vector retained. A kernel of dimension at least
two would leave a projective family and would require further work;
it is never replaced by sampled lines. In each of the ten cases in
the table, the complete enumeration below encounters no such family
before all branches are excluded.

The exclusion criterion is monotone. If a partial extension
$U+F_0\subset E$ already has an exterior-map rank larger than a bound
in (P17.1), every completion has at least that rank and is impossible.
We may therefore discard that branch immediately. Otherwise we
enumerate all its next invariant lines. This proves by induction that
every possible $E$ remains in the enumeration until an actual rank
contradiction excludes it.

Here is an integral specification of every matrix used. A basis
vector of $A$ is indexed by $Q\in S$ and a variable basis vector by
$v$. Its weight in $A\otimes B$ is
$\operatorname{wt}(v)-\operatorname{wt}(Q)$. In each weight containing
one image vector from $U$, eliminate the last coordinate in
lexicographic $(Q,v)$ order. Its coefficient is $1$ or $-1$. The
remaining coordinate vectors form the fixed quotient basis, and the
projection and inclusion matrices satisfy
\[
 QU=0,\qquad QX=1.
\]
For an adjacent transvection on variables let $N$ be its matrix on
$B$, and let $R$ be its induced matrix on the quadratic space $W$.
Replacing an index of a two-subset gives the ordinary exterior sign;
the order-ideal property ensures that the replacement stays in $S$.
The induced matrix on $A\otimes B$ is
\[
 G=-R^{\mathsf t}\otimes1+1\otimes N.
\]
The checker verifies the integer equivariance identity
$GU=-UN^{\mathsf t}$ and forms the quotient operator $QGX$.
It then forms each successive quotient by exact rational row
reduction, computes the common kernels in every weight, and retains
all distinct resulting subspaces. These are computations over
$\mathbb Q$, valid after every characteristic-zero field extension.

For the listed cases set $r=17$ and $\dim A=10$. The successive
quotient dimensions required are $17-m$. The numbers of distinct
spaces generated at each depth, before rank pruning, are
\[
\begin{array}{c|l|l}
 \text{case}&\text{generated at successive depths}
             &\text{discarded at successive depths}\\\hline
 1&3&3\\
 2&5,18&0,18\\
 3&5,19&0,19\\
 4&6,25&0,25\\
 5&6,25,84&0,0,84\\
 6&6,27,97&0,0,97\\
 7&5,19,57&0,4,57\\
 8&5,21,73&0,1,73\\
 9&4,15,36&0,5,36\\
 10&6,25,85,249&0,0,14,249
\end{array}
\]
Every last layer is empty after pruning. The self-contained
\path{verify_det10_borel.py} reconstructs all tensors, quotient
matrices, root operators, and flag successors directly from the ten
displayed subsets. It imports no project data, search routine, or
symbolic rank oracle. Its exact rational row reduction checks
$22120$ weight kernels; none requiring a projective choice is
encountered. For each of the $685$ rank exclusions it selects and
independently evaluates a nonzero exterior-map minor modulo the
verified prime $1000003$, of order
\[
 9\cdot17+1=154\quad\hbox{or}\quad(m-1)17+1.
\]
It writes the complete flag spaces, minor indices, and nonzero residues
to \path{det10_borel_certificate.json}. The modular products are less
than $2^{40}$ in magnitude and are held in signed 64-bit integers;
all quotient arithmetic uses arbitrary-precision rational numbers.
The residues certify nonzero integer minors, hence the stated
characteristic-zero rank contradictions. There is no admissible
seventeen-plane $E$ in any case. Transposition gives the other eight
subsets, and the finding follows.
\end{proof}

\paragraph{Dictionary and remaining gate.}
The argument takes place on a proper Grassmannian incidence with
determinantal rank conditions. Its fixed locus is refined by relative
flag incidences for the quotient of the tautological bundle. The
split torus gives the character summands, and the adjacent root
subgroups supply the explicit equivariant endomorphisms used to form
their common kernels. The resulting kernels and their projective
spaces are genuine bundle loci; their dimensions are not inferred
from representation classes in $K_0$. These are the same EGA~II,
SGA~3 and FGA constructions located in the common dictionary.
Together with the earlier six Koszul cases, P-17 reduces the
determinant ten-plane test for a bound forty-four to the remaining
$35$ Borel ideals. The additional exploratory files record smaller
extension-space gates for them. Those exploratory survivors do not
constitute a border-rank upper bound or a proof of the uniform lower
bound eighteen.
